\chapter*{Conclusion}
\addcontentsline{toc}{chapter}{Conclusion}  
\markboth{Conclusion}{Conclusion}

\paragraph{Chapter 1}
In the first chapter, I presented the physical concepts underlying the cooling scheme of our experiment, 
providing the elements necessary for a new PhD student -- or a similar cold-atom experiment -- to operate 
the setup and understand the physics behind it. Building on the previous theses from our team, this work 
also presents a new optical setup for filtering the MOGlabs ECDL laser used for the photoassociation (PA) 
process. The purpose of this filtering cavity is to suppress spontaneous emission on far-detuned states, 
which otherwise arises from the broadening of the laser spectrum due to amplified spontaneous emission.

\paragraph{Chapter 2}
The second chapter presents the results of our team's article published in PRX Quantum during my PhD. This 
work describes the quantum platform we engineered to measure physical properties of our Fermi gas using 
quantum-enhanced protocols. The article begins with the characterization of the pseudo-spin qubits we 
engineered: their coherence time reaches more than three seconds, owing to the low sensitivity of 
$^{87}$Sr's ground state to external perturbations.

The robustness of these pseudo-spin qubits then allowed us to extend the scheme to a multi-state 
interferometer -- made possible by the large Hilbert space accessible through the nuclear spin $I = 9/2$ 
of $^{87}$Sr -- giving access to physical quantities such as the linear and quadratic components of the 
light shift induced by the two-photon Raman process used for spin manipulation.

Finally, we demonstrate a measurement scheme for the non-commuting observables $\hat{s}^x, \hat{s}^y, 
\hat{s}^z$, which we show is equivalent to the production of Positive-Operator-Valued Measures (POVMs); 
this equivalence is derived in a dedicated section following the included article.

\paragraph{Chapter 3}
The last chapter presents our results on the photoassociation of $^{87}$Sr, a process that had remained 
unexplored prior to this work. The chapter opens with the molecular-physics background necessary to 
interpret the spectroscopic measurements, complemented by the analysis of temperature and atom-number 
decay as a function of illumination time. We encountered numerous technical difficulties along the way, 
which are documented here for the reader's reference and should be kept in mind more generally when 
performing this type of spectroscopy.

The observed lines display a surprisingly rich structure in both their strength and linewidth. The most 
plausible interpretation is a predissociation mechanism: free atoms are excited to a vibrational state 
close in energy to a coupled atomic potential, so that instead of forming a bound photoassociated 
molecule, they decay into an atomic kinetic continuum.
\paragraph{Outcome}
Several improvements to the experiment are already underway. To improve the resolution of the cloud imaging, 
a microscope has been built by a member of the team and is ready to be implemented; this system would also 
give access to the atom number before time-of-flight, providing a more accurate estimate at every step of the 
sequence. In parallel, a full replacement for the current transverse cooling system has been constructed and 
is ready to be installed on the main experiment table, which will consequently increase the number of atoms 
delivered to the vacuum cell.

On the scientific side, the interferometry scheme developed in this thesis also suggests a way to probe the 
SU(N) symmetry of the fermionic species \hl{compléter : lien précis avec la mesure des interactions}, opening 
a possible new diagnostic tool for interactions in this system.

Understanding the physical origin of the line broadening -- expected from the ongoing optical pumping studies 
-- is the next necessary step: it should make a larger fraction of the $^{87}$Sr spectroscopy accessible, and 
in turn allow the observed molecular lines to be unambiguously assigned to their associated molecular 
potentials. This identification is the last missing piece toward the original goal of this work: driving the 
system, through spin-selective photoassociative losses, into a highly-entangled Dicke state with well-defined 
spin components.